\begin{figure}[t]
\centering
\begin{tikzpicture}[
    font=\scriptsize,
    node distance=0.0cm,
    every node/.style={
        rotate=90,
        minimum width=3.3cm,
        minimum height=0.7cm,
        text centered,
        draw,
        fill=gray!10,
        rounded corners,
        inner xsep=2pt,
        inner ysep=6pt
    }
]

\node (stubCall)  [fill=blue!10]    {Call Stub 1};
\node (payload)   [below=of stubCall, fill=green!10]  {Wrapped Inst. / payload};
\node (pad1)      [below=of payload]                  {PAD (NOP)};
\node (stub2)     [below=of pad1, fill=yellow!20]     {Stub 2};
\node (pad2)      [below=of stub2]                    {PAD (NOP)};
\node (stub1)     [below=of pad2, fill=blue!10]       {Stub 1};
\node (pad3)      [below=of stub1]                    {PAD (NOP)};
\node (payloadPtr)[below=of pad3, fill=red!20]        {Pointer to Payload};
\node (pad4)      [below=of payloadPtr]               {PAD (NOP)};
\node (memZone)   [below=of pad4, fill=purple!10]     {Memory Zone};

\end{tikzpicture}
\Description{A schematic separates the initial call, instruction wrappers, redirector components, padding, offset information, and memory region.}
\caption{PII payload layout. Adapted from Fig.~1 of~\cite{zone_secrypt_2026}.}
\label{fig:pii_layout}
\end{figure}